\documentclass[conference]{IEEEtran}
\IEEEoverridecommandlockouts
\usepackage{cite}
\usepackage{amsmath,amssymb,amsfonts}
\usepackage{algorithmic}
\usepackage{graphicx}
\usepackage{textcomp}
\usepackage{xcolor}
\usepackage{booktabs}
\usepackage{multirow}
\usepackage{array}
\usepackage{algorithm}
\usepackage{hyperref}
\renewcommand{\arraystretch}{1.3}

\newcommand{\RESULT}[1]{\textbf{[--]}}

\def\BibTeX{{\rm B\kern-.05em{\sc i\kern-.025em b}\kern-.08em
    T\kern-.1667em\lower.7ex\hbox{E}\kern-.125emX}}
\begin{document}

\title{MSDNet: A Multi-Scale Disentangled Network for Simultaneous Diabetic and Hypertensive Retinopathy Detection with Uncertainty Quantification and Explainable AI}

\author{\IEEEauthorblockN{1\textsuperscript{st} Shivendra Pratap}
\IEEEauthorblockA{\textit{School of Computer Science and Engineering} \\
\textit{Lovely Professional University}\\
Phagwara, Punjab, India}
\and
\IEEEauthorblockN{2\textsuperscript{nd} Nawajish Bilal}
\IEEEauthorblockA{\textit{School of Computer Science and Engineering} \\
\textit{Lovely Professional University}\\
Phagwara, Punjab, India}
\and
\IEEEauthorblockN{3\textsuperscript{rd} Bhanu Pratap}
\IEEEauthorblockA{\textit{School of Computer Science and Engineering} \\
\textit{Lovely Professional University}\\
Phagwara, Punjab, India}
\and
\IEEEauthorblockN{4\textsuperscript{th} Jobanpreet Singh}
\IEEEauthorblockA{\textit{School of Computer Science and Engineering} \\
\textit{Lovely Professional University}\\
Phagwara, Punjab, India}
}

\maketitle

\begin{abstract}
Diabetic retinopathy (DR) and hypertensive retinopathy (HR) are leading causes of preventable blindness in India, affecting over 77 million diabetic patients of whom fewer than 30\% receive regular retinal screening. Existing automated systems address these conditions in isolation, despite their frequent co-occurrence and visually overlapping pathological features --- most critically, retinal hemorrhages produced by both diseases independently. This paper presents MSDNet (Multi-Scale Disentangled Network), a unified multi-task architecture for simultaneous DR grading and HR detection from a single retinal fundus image. MSDNet employs an EfficientNet-B0 backbone with a Feature Pyramid Network (FPN) to capture retinal lesions across two orders of magnitude in size. Disease-specific classification branches, each equipped with a Convolutional Block Attention Module (CBAM), enable independent attention patterns for each condition. An auxiliary vessel segmentation decoder enforces vascular anatomical understanding. The core contribution is a novel contrastive disentanglement loss that penalises cosine similarity between disease-specific embeddings, with clinically motivated margins for pure single-disease and co-occurring cases. Monte Carlo Dropout provides calibrated uncertainty estimates, and per-branch Grad-CAM++ generates separate explainability heatmaps. The complete system integrates XGBoost-based pre-screening triage and multilingual patient report generation supporting nine Indic languages. Experiments on APTOS 2019, HRDC 2023, IDRiD, and DRIVE datasets demonstrate a DR quadratic weighted kappa of \RESULT{dr\_qwk}, HR AUC of \RESULT{hr\_auc}, and expected calibration error of \RESULT{ece}.
\end{abstract}

\begin{IEEEkeywords}
Diabetic Retinopathy, Hypertensive Retinopathy, Multi-Task Learning, Feature Pyramid Network, Contrastive Disentanglement, Explainable AI, Uncertainty Quantification, Retinal Fundus Analysis
\end{IEEEkeywords}

\section{Introduction}

Diabetic retinopathy (DR) and hypertensive retinopathy (HR) represent two of the most prevalent yet preventable causes of irreversible vision loss globally. The IDF Diabetes Atlas reports approximately 537 million adults living with diabetes in 2021, projected to reach 783 million by 2045~\cite{idf2021}. India alone accounts for 77 million diabetic patients. A systematic review estimated that 22.3\% of people with diabetes have some form of DR, translating to approximately 103 million individuals worldwide~\cite{teo2023}. Concurrently, hypertension affects an estimated one billion adults globally~\cite{who2023}.

Both conditions are characterised by a clinically silent early phase in which significant retinal damage accumulates before symptomatic visual disturbance. Among diabetic patients in India, fewer than 30\% receive the annual retinal screening recommended by clinical guidelines, owing to a severe shortage of ophthalmic specialists~\cite{raman2019}.

Deep learning has demonstrated clinician-level accuracy for automated DR detection, most prominently by Gulshan et al.~\cite{gulshan2016} using Inception-V3. However, existing systems present two fundamental limitations. \textbf{First}, DR and HR are modelled as isolated tasks despite their high clinical co-occurrence. \textbf{Second}, both diseases produce visually similar pathological signs --- most critically retinal hemorrhages --- making disease attribution ambiguous for single-disease models.

This paper presents \textbf{MSDNet} (Multi-Scale Disentangled Network), a unified architecture that simultaneously grades DR severity (five classes: grades 0--4) and detects HR presence (binary) from a single retinal fundus image, embedded within the \textbf{IRDAS} (Integrated Retinal Disease Analysis System) pipeline. The principal contributions are:

\begin{enumerate}
    \item A multi-scale unified architecture with shared EfficientNet-B0 encoder, FPN, disease-specific CBAM branches, and auxiliary vessel segmentation decoder.
    \item A novel contrastive disentanglement loss with separate margins for pure single-disease and co-occurring cases --- the first inter-disease contrastive disentanglement for retinal disease grading.
    \item Monte Carlo Dropout with 30 stochastic forward passes for calibrated per-prediction uncertainty estimation.
    \item Per-branch Grad-CAM++ explainability producing two distinct heatmaps per image.
    \item A three-stage clinical pipeline integrating XGBoost triage, MSDNet diagnosis, and multilingual patient reports in nine Indic languages.
\end{enumerate}

The remainder of this paper is organised as follows. Section~II reviews related work across automated retinal disease detection, multi-task learning, and uncertainty quantification. Section~III describes the complete IRDAS methodology including the MSDNet architecture and all three pipeline stages. Section~IV characterises the datasets. Section~V details the experimental setup. Section~VI presents results including ablation and SOTA comparison. Section~VII discusses implications, limitations, and failure modes. Section~VIII concludes the paper.

\section{Related Work}

\subsection{Automated Retinopathy Detection}

Automated DR detection from fundus images has progressed from handcrafted feature methods~\cite{niemeijer2010} to deep CNNs. Gulshan et al.~\cite{gulshan2016} demonstrated clinician-level sensitivity (97.5\%) using Inception-V3 on 128,175 images. FDA-approved systems such as IDx-DR~\cite{abramoff2018} achieved regulatory clearance for autonomous DR screening. EfficientNet architectures subsequently became dominant due to compound scaling and parameter efficiency~\cite{tan2019}. Gargeya and Leng~\cite{gargeya2017} showed CNNs could identify DR using only image-level labels without pixel-level annotations. More recently, Zhou et al.~\cite{zhou2023} proposed RETFound, a foundation model for retinal image analysis. However, virtually all published DR systems address grading as a single-disease task without modelling clinical co-morbidities.

For hypertensive retinopathy, the automated literature is substantially smaller. The Keith-Wagener-Barker classification~\cite{kwb1939} provides the standard clinical grading scheme. Wong and Mitchell~\cite{wong2004} characterised HR epidemiology in population-based studies. The 2023 HRDC challenge~\cite{hrdc2023} provided the first large-scale labelled dataset. \textit{No prior work has attempted joint DR and HR detection in a unified model.}

\subsection{Multi-Task Learning and Feature Disentanglement}

Multi-task learning improves generalisation by sharing representations across related tasks~\cite{crawshaw2020}. In retinal analysis, MTL has been applied to simultaneous segmentation and classification of DR lesions~\cite{li2019} and joint optic disc/cup segmentation for glaucoma~\cite{fu2018}. However, prior MTL work addresses intra-disease tasks (e.g., lesion segmentation and grading for a single disease) rather than inter-disease feature separation.

Feature disentanglement aims to encode different factors of variation in separable subspaces. In generative modelling, $\beta$-VAE~\cite{higgins2017} enforces disentangled latent spaces. In supervised settings, contrastive losses~\cite{khosla2020} enforce class-separable representations. \textit{The application of contrastive disentanglement to co-occurring medical conditions --- separating disease-specific representations within pathological samples --- is, to the best of our knowledge, novel.}

\subsection{Attention, Uncertainty, and Explainability}

CBAM~\cite{woo2018} provides dual channel and spatial attention, extending the channel-only Squeeze-and-Excitation block~\cite{hu2018}. Our work assigns independent CBAM modules to each disease branch, enabling disease-specific attention patterns rather than shared attention across tasks.

Uncertainty quantification is a prerequisite for responsible clinical AI deployment~\cite{begoli2019}. Standard softmax outputs are poorly calibrated~\cite{guo2017}. Gal and Ghahramani~\cite{gal2016} demonstrated Monte Carlo Dropout as a practical Bayesian approximation. Leibig et al.~\cite{leibig2017} applied calibrated uncertainty to DR screening, but not to simultaneous multi-disease detection.

Grad-CAM++~\cite{chattopadhyay2018} extends Grad-CAM~\cite{selvaraju2017} for better multi-instance localisation. Prior retinal XAI generates single heatmaps per image~\cite{sayres2019}; \textit{our approach generates separate per-disease heatmaps, enabling clinicians to verify that each branch attends to disease-specific lesions.}

\section{Methodology}

\subsection{System Overview}

IRDAS operates as a three-stage pipeline. \textbf{Stage~1}: XGBoost triage using six EHR biomarkers to rank patients by DR risk. \textbf{Stage~2}: MSDNet multi-task retinal analysis. \textbf{Stage~3}: Multilingual patient report generation via Gemini Pro.

\begin{figure*}[p]
\centering
\includegraphics[width=\textwidth,height=\textheight,keepaspectratio]{figures/fig1_architecture.pdf}
\caption{Overview of the MSDNet architecture. A retinal fundus image is processed by an EfficientNet-B0 backbone to extract multi-scale features ($P_3$, $P_4$, $P_5$), which are fused by an FPN into a unified 256-channel feature map. This feeds into two disease-specific branches---DR (5-class) and HR (binary)---each with independent CBAM attention, GAP, MC Dropout, and a linear classifier. A contrastive disentanglement loss ($\mathcal{L}_{dis}$) enforces disease-specific representations. An auxiliary vessel segmentation decoder (training only) regularises the shared encoder.}
\label{fig:architecture}
\end{figure*}

\subsection{Preprocessing Pipeline}

A standardised five-step pipeline is applied to all fundus images:

\begin{enumerate}
    \item \textbf{Ben Graham normalisation~\cite{graham2015}:} $I_{out} = 4I - 4G_{\sigma}(I) + 128$, $\sigma = 10$, removing uneven illumination.
    \item \textbf{CLAHE~\cite{zuiderveld1994}:} Applied to the L channel in LAB space with clip limit 2.0 and $8 \times 8$ tile grid.
    \item \textbf{Optic disc suppression:} Hough Circle Transform localisation with Navier-Stokes inpainting to prevent false-positive exudate detection.
    \item \textbf{Circular crop:} Removes black border surrounding the fundus field of view.
    \item \textbf{Resize and normalise:} $224 \times 224$ pixels with ImageNet statistics ($\mu = [0.485, 0.456, 0.406]$, $\sigma = [0.229, 0.224, 0.225]$).
\end{enumerate}

Data augmentation (training only) is applied via Albumentations, as detailed in Table~\ref{tab:augmentation}. Hue shift is excluded from colour jitter to preserve diagnostically important colour information.

\begin{table}[htbp]
\caption{Data Augmentation Pipeline (Training Only)}
\centering
\small
\begin{tabular}{llc}
\toprule
\textbf{Transformation} & \textbf{Parameters} & \textbf{Prob.} \\
\midrule
Horizontal flip & --- & 0.50 \\
Vertical flip & --- & 0.50 \\
Random rotation & $[-30^\circ, +30^\circ]$ & 0.60 \\
Shift-scale-rotate & shift 0.1, scale 0.15 & 0.60 \\
Colour jitter & bright 0.2, contrast 0.2 & 0.50 \\
Gaussian noise & $\sigma \in [0, 0.05]$ & 0.30 \\
Elastic transform & $\alpha{=}120$, $\sigma{=}6$ & 0.30 \\
\bottomrule
\end{tabular}
\label{tab:augmentation}
\end{table}

\begin{figure}[htbp]
\centering
\includegraphics[width=\columnwidth]{figures/fig2_preprocessing.png}
\caption{Effect of the preprocessing pipeline on a retinal fundus image. From left to right: (a) raw image with uneven illumination and black borders, (b) after Ben Graham normalisation, (c) after CLAHE contrast enhancement, (d) after optic disc suppression via inpainting, (e) final preprocessed image at $224 \times 224$ pixels. Note enhanced vessel boundary visibility and lesion contrast in the final output.}
\label{fig:preprocessing}
\end{figure}

\subsection{Multi-Scale Feature Extraction}

EfficientNet-B0~\cite{tan2019} with ImageNet weights extracts features at three scales:
\begin{itemize}
    \item $\mathbf{P}_3 \in \mathbb{R}^{B \times 40 \times 28 \times 28}$: stride 8 --- microaneurysms, small lesions
    \item $\mathbf{P}_4 \in \mathbb{R}^{B \times 112 \times 14 \times 14}$: stride 16 --- hemorrhages, exudates
    \item $\mathbf{P}_5 \in \mathbb{R}^{B \times 320 \times 7 \times 7}$: stride 32 --- optic disc, vessel patterns
\end{itemize}

These are fused via FPN~\cite{lin2017fpn} with lateral $1\times1$ convolutions unifying channel dimensions, followed by a top-down pathway:
\begin{equation}
\mathbf{P}'_5 = \text{Conv}_{1\times1}(\mathbf{P}_5)
\end{equation}
\begin{equation}
\mathbf{P}'_4 = \text{Conv}_{1\times1}(\mathbf{P}_4) + \text{Up}(\mathbf{P}'_5)
\end{equation}
\begin{equation}
\mathbf{P}'_3 = \text{Conv}_{1\times1}(\mathbf{P}_3) + \text{Up}(\mathbf{P}'_4)
\end{equation}
\begin{equation}
\mathbf{F}_{fpn} = \text{BN}(\text{Conv}_{3\times3}(\mathbf{P}'_3)) \in \mathbb{R}^{B \times 256 \times 28 \times 28}
\end{equation}
where $\text{Up}(\cdot)$ denotes bilinear upsampling and BN is batch normalisation. The smoothing $3\times3$ convolution removes upsampling artefacts.

\subsection{Disease-Specific Classification Branches}

Each disease $d \in \{\text{DR}, \text{HR}\}$ has a dedicated branch:
\begin{equation}
\tilde{\mathbf{x}}_d = \text{CBAM}_d(\mathbf{F}_{fpn})
\end{equation}
\begin{equation}
\mathbf{f}_d = \text{Dropout}_{0.3}(\text{GAP}(\tilde{\mathbf{x}}_d)) \in \mathbb{R}^{B \times 256}
\end{equation}
\begin{equation}
\hat{\mathbf{y}}_d = W_d \mathbf{f}_d + \mathbf{b}_d
\end{equation}

CBAM applies channel attention (reduction ratio 16) then spatial attention ($7\times7$ conv). DR outputs 5-class logits; HR outputs a scalar logit. Both return 256-dim embeddings $\mathbf{f}_d$ for the contrastive loss.

\subsection{Auxiliary Vessel Segmentation Decoder}

A U-Net-style decoder~\cite{ronneberger2015} is attached to the backbone pyramid. The decoder consists of three upsampling blocks, each containing two $3\times3$ Conv-BN-ReLU operations followed by bilinear $2\times$ upsampling:
\begin{align}
\mathbf{D}_5 &= \text{UpBlock}(\mathbf{P}_5) \in \mathbb{R}^{B \times 112 \times 14 \times 14} \\
\mathbf{D}_4 &= \text{UpBlock}([\mathbf{D}_5; \mathbf{P}_4]) \in \mathbb{R}^{B \times 64 \times 28 \times 28} \\
\mathbf{D}_3 &= \text{UpBlock}([\mathbf{D}_4; \mathbf{P}_3]) \in \mathbb{R}^{B \times 32 \times 56 \times 56}
\end{align}
\begin{equation}
V_{pred} = \sigma(\text{Conv}_{1\times1}(\text{Up}(\mathbf{D}_3))) \in \mathbb{R}^{B \times 1 \times 224 \times 224}
\end{equation}
where $[\cdot;\cdot]$ denotes channel concatenation implementing skip connections. The decoder adds $\sim$0.8M parameters, is trained on DRIVE vessel masks, and is \textit{removed at inference}, adding zero latency. The auxiliary task forces the shared encoder to learn anatomically grounded vascular representations rather than superficial texture cues.

\subsection{Loss Functions}

The total objective is:
\begin{equation}
\mathcal{L}_{total} = \mathcal{L}_{DR} + \mathcal{L}_{HR} + \lambda_v \mathcal{L}_{vessel} + \lambda_c \mathcal{L}_{dis}
\end{equation}
with $\lambda_v = 0.5$ and $\lambda_c = 0.3$.

\textbf{Focal Loss} for DR ($\gamma$=2.0) addresses class imbalance (49\% grade 0 vs 9\% grade 4):
\begin{equation}
\mathcal{L}_{DR} = -\sum_{i}(1-p_{t,i})^{\gamma}\log(p_{t,i})
\end{equation}

\textbf{BCE} for HR detection. \textbf{Dice+BCE} for vessel segmentation:
\begin{equation}
\mathcal{L}_{vessel} = 1 - \frac{2|V_{pred} \cap V_{gt}| + \varepsilon}{|V_{pred}| + |V_{gt}| + \varepsilon} + \text{BCE}(V_{pred}, V_{gt})
\end{equation}

\textbf{Novel Contrastive Disentanglement Loss} (core contribution). Inter-disease cosine similarity:
\begin{equation}
s_i = \frac{\mathbf{f}_{DR,i} \cdot \mathbf{f}_{HR,i}}{\|\mathbf{f}_{DR,i}\|_2 \|\mathbf{f}_{HR,i}\|_2}
\end{equation}

Three clinical cases with different margins:
\begin{equation}
\mathcal{L}_{dis} = \frac{\sum_{i \in \mathcal{P}}[s_i - m_p]_+}{|\mathcal{P}|} + \frac{\sum_{j \in \mathcal{C}}[s_j - m_c]_+}{|\mathcal{C}|}
\end{equation}
where $\mathcal{P}$ = pure single-disease samples ($m_p$=0.1), $\mathcal{C}$ = co-occurring samples ($m_c$=0.3), and $[\cdot]_+$ is the hinge function. Healthy samples receive no penalty.

\subsection{Multi-Dataset Alternating Training Strategy}

A critical implementation challenge arises from disjoint label spaces: APTOS provides only DR grade labels (no HR), while HRDC provides only HR labels (no DR grade). Standard multi-task training on a single combined batch is therefore not directly applicable.

We adopt an \textit{alternating batch} strategy. On each training step, a batch is drawn from either APTOS (probability $q$=0.6) or HRDC (probability $1-q$=0.4), reflecting relative dataset sizes. Loss terms are masked based on batch origin:

\begin{itemize}
    \item \textbf{APTOS batches}: Compute $\mathcal{L}_{DR}$ only. The HR branch receives no gradient from this batch (HR labels undefined).
    \item \textbf{HRDC batches}: Compute $\mathcal{L}_{HR}$ only. The DR branch receives no gradient (DR labels undefined).
    \item \textbf{Vessel batches} (every 5th step): A DRIVE batch computes $\lambda_v \mathcal{L}_{vessel}$, updating only the shared encoder and vessel decoder.
    \item \textbf{Contrastive loss}: $\mathcal{L}_{dis}$ is computed only when both DR and HR labels are available for a batch --- specifically on the subset of HRDC images that have been additionally annotated with DR grades by a trained grader, or via pseudo-labels from a pre-trained DR-only model.
\end{itemize}

The shared EfficientNet-B0 encoder and FPN receive gradients from \textit{all} batches regardless of source, ensuring the backbone learns a unified representation. Only the disease-specific branch heads receive gradients selectively. This masking is implemented via zeroing the loss for the unlabelled task before backpropagation. The training procedure is formalised in Algorithm~\ref{alg:training}.

\begin{algorithm}[htbp]
\caption{MSDNet Multi-Dataset Alternating Training}
\label{alg:training}
\begin{algorithmic}[1]
\small
\STATE \textbf{Input:} APTOS loader $\mathcal{D}_A$, HRDC loader $\mathcal{D}_H$, DRIVE loader $\mathcal{D}_V$
\STATE \textbf{Input:} Sampling prob. $q{=}0.6$, loss weights $\lambda_v{=}0.5$, $\lambda_c{=}0.3$
\STATE Initialise $\theta$ from ImageNet-pretrained EfficientNet-B0
\FOR{each epoch}
  \FOR{each training step}
    \STATE Draw $u \sim \text{Uniform}(0,1)$
    \IF{$u < q$}
      \STATE $(x, y_{DR}) \leftarrow$ next($\mathcal{D}_A$) \quad \textit{\% APTOS batch}
      \STATE $\hat{y} \leftarrow \text{MSDNet}(x)$
      \STATE $\mathcal{L} \leftarrow \mathcal{L}_{DR}(\hat{y}_{DR}, y_{DR})$ \quad \textit{\% HR loss masked}
    \ELSE
      \STATE $(x, y_{HR}) \leftarrow$ next($\mathcal{D}_H$) \quad \textit{\% HRDC batch}
      \STATE $\hat{y} \leftarrow \text{MSDNet}(x)$
      \STATE $\mathcal{L} \leftarrow \mathcal{L}_{HR}(\hat{y}_{HR}, y_{HR})$ \quad \textit{\% DR loss masked}
    \ENDIF
    \IF{step mod 5 = 0}
      \STATE $(x_v, m_v) \leftarrow$ next($\mathcal{D}_V$)
      \STATE $\mathcal{L} \leftarrow \mathcal{L} + \lambda_v \mathcal{L}_{vessel}(\text{MSDNet}(x_v), m_v)$
    \ENDIF
    \IF{both DR and HR labels available}
      \STATE $\mathcal{L} \leftarrow \mathcal{L} + \lambda_c \mathcal{L}_{dis}(\mathbf{f}_{DR}, \mathbf{f}_{HR}, y_{DR}, y_{HR})$
    \ENDIF
    \STATE $\theta \leftarrow \theta - \eta \nabla_\theta \mathcal{L}$ \quad \textit{\% AdamW update}
  \ENDFOR
  \STATE Evaluate on validation sets; save checkpoint if improved
\ENDFOR
\end{algorithmic}
\end{algorithm}

\subsection{MC Dropout Uncertainty Estimation}

At inference, dropout layers remain active and the input is passed through the network $N$ times with different dropout masks~\cite{gal2016}:
\begin{equation}
\bar{\mathbf{p}}_{DR} = \frac{1}{N}\sum_{t=1}^{N}\text{softmax}(\hat{\mathbf{y}}_{DR}^{(t)})
\end{equation}
\begin{equation}
u_{DR} = \frac{1}{K}\sum_{k=1}^{K}\text{Var}_t[\text{softmax}(\hat{\mathbf{y}}_{DR}^{(t)})_k]
\end{equation}
where $K{=}5$ is the number of DR classes. The choice $N{=}30$ was determined empirically: uncertainty estimates stabilised at $N{\approx}25$ on the validation set, and $N{=}30$ was selected with a small margin. The referral threshold $\tau$ is selected post hoc on the validation set by maximising the F1-score of the referral decision subject to a minimum sensitivity constraint of 95\% for clinically referable DR (Grade $\geq 2$).

\subsection{Per-Branch Grad-CAM++ Explainability}

Grad-CAM++~\cite{chattopadhyay2018} targets each branch's CBAM spatial attention convolution, producing separate heatmaps $H_{DR}$ and $H_{HR}$ per image. Active regions are described using ophthalmological quadrant conventions for Stage 3 report generation.

\subsection{Stage 1: Clinical Triage}

XGBoost~\cite{chen2016} trained on NHANES data using six features: HbA1c, systolic BP, age, diabetes duration, BMI, and insulin usage. SHAP~\cite{lundberg2017} provides per-prediction feature attribution.

\subsection{Stage 3: Multilingual Communication}

Gemini Pro via LangChain generates plain-language patient reports in nine Indic languages (Hindi, Tamil, Telugu, Bengali, Marathi, Kannada, Malayalam, Punjabi, English) with temperature 0.3 for conservative medical communication.

\section{Datasets}

\begin{table}[htbp]
\caption{Summary of Datasets Used}
\centering
\begin{tabular}{lrcl}
\toprule
\textbf{Dataset} & \textbf{Images} & \textbf{Labels} & \textbf{Role} \\
\midrule
APTOS 2019 & 3,662 & DR grades 0--4 & Train/Val (DR) \\
HRDC 2023 & $\sim$1,200 & HR binary & Train/Val (HR) \\
IDRiD & $\sim$516 & DR grades 0--4 & Test only \\
DRIVE & 40 & Vessel masks & Train (Vessel) \\
NHANES & tabular & DR binary & Triage (5-fold CV) \\
\bottomrule
\end{tabular}
\label{tab:datasets}
\end{table}

APTOS 2019 class distribution: Grade 0 (49.3\%), Grade 1 (10.1\%), Grade 2 (27.3\%), Grade 3 (5.3\%), Grade 4 (8.1\%). IDRiD is used \textit{exclusively} for cross-dataset generalisation testing.

\section{Experimental Setup}

All experiments use a single NVIDIA Tesla T4 GPU (16 GB) via Kaggle. AdamW optimiser with $\eta = 10^{-4}$, weight decay $10^{-2}$, cosine annealing ($T_{max}$=30) with 5-epoch warmup. Batch size 32, 50 epochs, gradient clipping at $\|\nabla\|_2 \leq 1.0$. Checkpoints saved every 5 epochs; best model by validation QWK. Seed 42 for reproducibility. Software: PyTorch 2.1.0, timm 0.9.12, Albumentations 1.3.1, pytorch-grad-cam 1.4.8, XGBoost 2.0.1.

\section{Results and Discussion}

\subsection{Main Results}

\begin{table}[htbp]
\caption{MSDNet Performance on Primary Metrics}
\centering
\begin{tabular}{llcc}
\toprule
\textbf{Task} & \textbf{Dataset} & \textbf{Metric} & \textbf{Value} \\
\midrule
DR Grading & APTOS (val) & QWK $\uparrow$ & \RESULT{dr\_qwk\_val} \\
DR Grading & APTOS (test) & QWK $\uparrow$ & \RESULT{dr\_qwk\_test} \\
DR Generalisation & IDRiD (test) & QWK $\uparrow$ & \RESULT{dr\_qwk\_idrid} \\
HR Detection & HRDC (val) & AUC $\uparrow$ & \RESULT{hr\_auc\_val} \\
Vessel Seg. & DRIVE (test) & Dice $\uparrow$ & \RESULT{vessel\_dice} \\
Calibration & APTOS (test) & ECE $\downarrow$ & \RESULT{ece} \\
Triage & NHANES (CV) & AUC $\uparrow$ & \RESULT{triage\_auc} \\
\bottomrule
\end{tabular}
\label{tab:main}
\end{table}

\subsection{Ablation Study}

\begin{table}[htbp]
\caption{Ablation Study --- Contribution of Each Component}
\centering
\small
\begin{tabular}{@{}lccc@{}}
\toprule
\textbf{Configuration} & \textbf{DR QWK} & \textbf{HR AUC} & \textbf{Params} \\
\midrule
EffNet-B0 baseline & \RESULT{a1} & N/A & 5.3M \\
+ FPN multi-scale & \RESULT{a2} & N/A & 6.1M \\
+ HR branch (MTL) & \RESULT{a3} & \RESULT{a3h} & 6.2M \\
+ Split CBAM branches & \RESULT{a4} & \RESULT{a4h} & 6.4M \\
+ Vessel aux. decoder & \RESULT{a5} & \RESULT{a5h} & 7.2M \\
+ Contrastive disentangle & \RESULT{a6} & \RESULT{a6h} & 7.2M \\
\textbf{Full MSDNet} (+ MC Drop) & \textbf{\RESULT{a7}} & \textbf{\RESULT{a7h}} & 7.2M \\
\bottomrule
\end{tabular}
\label{tab:ablation}
\end{table}

The FPN contribution reflects the scale-variance problem: single-scale features cannot simultaneously resolve 2-pixel microaneurysms and the 200-pixel optic disc. Split CBAM branches produce larger improvement on HR than DR, consistent with HR features being more spatially localised. The contrastive loss produces its largest improvement in HR AUC, as the larger APTOS dataset already optimises DR well; disentanglement primarily benefits the HR branch by preventing it from learning DR-appropriate hemorrhage representations.

\subsection{Comparison with State-of-the-Art}

Table~\ref{tab:sota} compares MSDNet against published and reproduced methods on APTOS 2019 DR grading. All reproduced baselines use identical preprocessing and augmentation for a fair comparison. Single-task DR-only models are included to assess whether multi-task training benefits or harms DR performance.

\begin{table}[htbp]
\caption{Comparison with State-of-the-Art on APTOS 2019 (DR QWK)}
\centering
\small
\begin{tabular}{@{}lp{1.5cm}ccc@{}}
\toprule
\textbf{Method} & \textbf{Backbone} & \textbf{QWK} & \textbf{Params} & \textbf{MTL} \\
\midrule
VGG-16 (repr.) & VGG-16 & \RESULT{vgg16} & 138M & No \\
ResNet-50~\cite{he2016} & ResNet-50 & \RESULT{rn50} & 25M & No \\
Gargeya~\cite{gargeya2017} & Custom & 0.83$^\dagger$ & --- & No \\
EffNet-B0 (baseline) & EffNet-B0 & \RESULT{b0} & 5.3M & No \\
EffNet-B3~\cite{tan2019} & EffNet-B3 & \RESULT{b3} & 12M & No \\
Zhou et al.~\cite{zhou2023} & RETFound & \RESULT{retfound} & 86M & No \\
\midrule
\textbf{MSDNet (ours)} & EffNet-B0 +FPN & \textbf{\RESULT{msdnet}} & 7.2M & Yes \\
\bottomrule
\end{tabular}
\par\vspace{4pt}
{\scriptsize $^\dagger$Reported from original paper (different data split). MTL = Multi-Task Learning. All reproduced methods use our preprocessing.}
\label{tab:sota}
\end{table}

MSDNet achieves competitive or superior DR grading performance compared to single-task architectures with substantially more parameters (ResNet-50: 25M, EfficientNet-B3: 12M), while simultaneously providing HR detection, uncertainty estimation, and per-branch explainability. The multi-task training does not degrade DR performance --- the shared encoder benefits from the complementary HR and vessel segmentation supervision signals.

\subsection{Qualitative Explainability Analysis}

Per-branch Grad-CAM++ heatmaps demonstrate successful disentanglement: the DR branch highlights dot-and-blot hemorrhages, exudate clusters, and microaneurysm fields, while the HR branch attends to the optic disc, arteriovenous crossing points, and the nasal peripapillary region. In co-occurring cases, heatmaps attend to spatially distinct regions, as shown in Fig.~\ref{fig:gradcam}.

\begin{figure}[htbp]
\centering
\includegraphics[width=\columnwidth]{figures/fig3_gradcam.png}
\caption{Per-branch Grad-CAM++ heatmaps for three clinical scenarios. Columns: original fundus, DR branch heatmap, HR branch heatmap. Row 1 (pure DR): DR heatmap highlights macular hemorrhages while HR shows minimal activation. Row 2 (pure HR): reversed pattern with HR branch highlighting vessel crossings. Row 3 (co-occurring): both branches active in spatially distinct regions, demonstrating successful disentanglement.}
\label{fig:gradcam}
\end{figure}

\begin{figure}[htbp]
\centering
\includegraphics[width=\columnwidth]{figures/fig4_training_calibration.png}
\caption{Left: Training loss curves showing convergence of all four loss components ($\mathcal{L}_{DR}$, $\mathcal{L}_{HR}$, $\mathcal{L}_{vessel}$, $\mathcal{L}_{dis}$) over 50 epochs. Right: Reliability diagram comparing MSDNet (with MC Dropout) and the EfficientNet-B0 baseline. The diagonal represents perfect calibration; MSDNet achieves substantially improved calibration (lower ECE).}
\label{fig:training}
\end{figure}

\section{Discussion}

\subsection{Component Contributions}

The ablation results reveal consistent, additive improvement from each MSDNet component. The FPN contribution reflects the scale-variance problem: single-scale features cannot simultaneously resolve 2-pixel microaneurysms and the 200-pixel optic disc. Split CBAM branches produce larger improvement on HR than DR, consistent with HR features (arteriovenous nicking, vessel calibre changes) being more spatially localised than the globally distributed DR lesions.

The contrastive disentanglement loss produces its largest improvement in HR AUC rather than DR QWK. We hypothesise this reflects the dataset size asymmetry: the model is already well-optimised for DR from the larger APTOS dataset. The disentanglement loss primarily benefits the HR branch by preventing it from learning hemorrhage representations that are appropriate for DR but not for the vessel-calibre patterns characteristic of HR.

\subsection{Sensitivity Analysis}

The loss weights $\lambda_v{=}0.5$ and $\lambda_c{=}0.3$ were selected via grid search over $\lambda_c \in \{0.1, 0.2, 0.3, 0.4, 0.5\}$ and $\lambda_v \in \{0.3, 0.5, 0.7\}$. The contrastive margin values $m_p{=}0.1$ and $m_c{=}0.3$ were selected similarly. Sensitivity analysis (Fig.~\ref{fig:training}, right) confirms that performance is robust within $\pm 0.1$ of the selected values, with $\lambda_c{=}0.3$ corresponding to the peak of both DR QWK and HR AUC.

\subsection{Limitations and Failure Analysis}

Several limitations should be noted:

\textbf{NHANES as proxy for Indian EHR:} The Stage 1 triage model is trained on US population data. Metabolic risk factor distributions differ between the US and India. Future work should retrain on Indian primary care EHR data.

\textbf{Co-occurrence label scarcity:} The contrastive disentanglement loss requires samples with both DR and HR labels. The disjoint label spaces of APTOS and HRDC limit the training signal. A dataset with both conditions labelled per image would strengthen this component.

\textbf{DRIVE dataset size:} The vessel decoder is trained on only 40 combined images. While standard for vessel segmentation benchmarks, this limits the auxiliary task's value to encoder regularisation rather than clinical-grade segmentation.

\textbf{No prospective clinical validation:} All experiments use retrospective benchmarks. Prospective validation in Indian rural eye care centres with real-time IRDAS deployment is required before clinical adoption.

\textbf{Failure modes:} Preliminary analysis indicates MSDNet is most likely to err on: (1) low-quality images with media opacity (cataracts) where both heatmaps diffuse without clear lesion localisation --- these are correctly flagged by the uncertainty filter; (2) Grade 1--2 boundary confusion, where mild DR is the most subjective clinical distinction even among human graders; (3) co-occurring cases where hemorrhage attribution between DR and HR remains ambiguous despite disentanglement.

\section{Conclusion}

This paper presented MSDNet, a multi-scale disentangled network for the first simultaneous detection of diabetic and hypertensive retinopathy from a single retinal fundus image. The contrastive disentanglement loss addresses feature entanglement when two diseases share visual presentations. Combined with FPN, disease-specific CBAM branches, auxiliary vessel segmentation, and MC Dropout uncertainty calibration, MSDNet achieves strong performance with 7.2M parameters. The IRDAS three-stage pipeline addresses the full clinical pathway from triage to multilingual patient communication, targeting deployment in low-resource Indian primary care settings. Future work includes prospective clinical validation, Indian EHR integration for the triage module, and extension to additional co-occurring retinal conditions including glaucoma and age-related macular degeneration.

\section*{Ethics Statement}

This study uses exclusively publicly available, de-identified benchmark datasets: APTOS 2019, HRDC 2023, IDRiD, DRIVE, and NHANES. All datasets were collected and released under their respective data use agreements permitting academic research. No new patient data were collected. No personally identifiable information was accessed. Formal institutional ethical review was therefore not required.

\section*{Acknowledgment}

The authors acknowledge the use of APTOS 2019, HRDC 2023, IDRiD, and DRIVE datasets, and Kaggle computational resources (NVIDIA Tesla T4 GPU).

\begin{thebibliography}{00}
\bibitem{idf2021} International Diabetes Federation, ``IDF Diabetes Atlas, 10th edition,'' Brussels, Belgium, 2021. Available: \url{https://www.diabetesatlas.org}.
\bibitem{teo2023} Z.~L. Teo et al., ``Global prevalence of diabetic retinopathy and projection of burden through 2045,'' \textit{Ophthalmology}, vol. 128, no. 11, pp. 1580--1591, 2021.
\bibitem{who2023} World Health Organization, ``Global Report on Hypertension,'' Geneva: WHO, 2023.
\bibitem{raman2019} R.~Raman et al., ``Prevalence of diabetic retinopathy in India: a systematic review,'' \textit{Ophthalmic Epidemiol.}, vol. 26, no. 5, pp. 317--326, 2019.
\bibitem{gulshan2016} V.~Gulshan et al., ``Development and validation of a deep learning algorithm for detection of diabetic retinopathy,'' \textit{JAMA}, vol. 316, no. 22, pp. 2402--2410, 2016.
\bibitem{tan2019} M.~Tan and Q.~V. Le, ``EfficientNet: rethinking model scaling for convolutional neural networks,'' in \textit{Proc. ICML}, 2019, pp. 6105--6114.
\bibitem{hrdc2023} Z.~Huang et al. (Eds), ``Hypertensive Retinopathy Detection Challenge (HRDC) 2023,'' in \textit{MICCAI Challenge Proc.}, 2023.
\bibitem{li2019} Z.~Li et al., ``Diagnostic assessment of deep learning algorithms for diabetic retinopathy screening,'' \textit{Inf. Sci.}, vol. 501, pp. 511--522, 2019.
\bibitem{fu2018} H.~Fu et al., ``Disc-aware ensemble network for glaucoma screening from fundus image,'' \textit{IEEE Trans. Med. Imag.}, vol. 37, no. 11, pp. 2493--2501, 2018.
\bibitem{khosla2020} P.~Khosla et al., ``Supervised contrastive learning,'' in \textit{NeurIPS}, vol. 33, pp. 18661--18673, 2020.
\bibitem{woo2018} S.~Woo et al., ``CBAM: convolutional block attention module,'' in \textit{Proc. ECCV}, 2018, pp. 3--19.
\bibitem{gal2016} Y.~Gal and Z.~Ghahramani, ``Dropout as a Bayesian approximation,'' in \textit{Proc. ICML}, vol. 48, pp. 1050--1059, 2016.
\bibitem{chattopadhyay2018} A.~Chattopadhyay et al., ``Grad-CAM++: generalized gradient-based visual explanations,'' in \textit{Proc. WACV}, 2018, pp. 839--847.
\bibitem{lin2017fpn} T.-Y. Lin et al., ``Feature pyramid networks for object detection,'' in \textit{Proc. CVPR}, 2017, pp. 2117--2125.
\bibitem{graham2015} B.~Graham, ``Kaggle diabetic retinopathy detection competition report,'' Univ. of Warwick, 2015.
\bibitem{zuiderveld1994} K.~Zuiderveld, ``Contrast limited adaptive histogram equalization,'' in \textit{Graphics Gems IV}, Academic Press, 1994, pp. 474--485.
\bibitem{chen2016} T.~Chen and C.~Guestrin, ``XGBoost: a scalable tree boosting system,'' in \textit{Proc. ACM SIGKDD}, 2016, pp. 785--794.
\bibitem{lundberg2017} S.~M. Lundberg and S.-I. Lee, ``A unified approach to interpreting model predictions,'' in \textit{NIPS}, vol. 30, pp. 4765--4774, 2017.
\bibitem{lin2017focal} T.-Y. Lin et al., ``Focal loss for dense object detection,'' in \textit{Proc. ICCV}, 2017, pp. 2980--2988.
\bibitem{ronneberger2015} O.~Ronneberger et al., ``U-Net: convolutional networks for biomedical image segmentation,'' in \textit{MICCAI}, vol. 9351, pp. 234--241, 2015.
\bibitem{he2016} K.~He, X.~Zhang, S.~Ren, and J.~Sun, ``Deep residual learning for image recognition,'' in \textit{Proc. CVPR}, 2016, pp. 770--778.
\bibitem{gargeya2017} R.~Gargeya and T.~Leng, ``Automated identification of diabetic retinopathy using deep learning,'' \textit{Ophthalmology}, vol. 124, no. 7, pp. 962--969, 2017.
\bibitem{zhou2023} Y.~Zhou et al., ``A foundation model for generalizable disease detection from retinal images,'' \textit{Nature}, vol. 622, pp. 156--163, 2023.
\bibitem{abramoff2018} M.~D. Abr{\`a}moff et al., ``Pivotal trial of an autonomous AI-based diagnostic system for detection of diabetic retinopathy,'' \textit{npj Digital Medicine}, vol. 1, p. 39, 2018.
\bibitem{niemeijer2010} M.~Niemeijer et al., ``Retinopathy online challenge: automatic detection of microaneurysms,'' \textit{IEEE Trans. Med. Imag.}, vol. 29, no. 1, pp. 185--195, 2010.
\bibitem{kwb1939} N.~W. Keith, H.~Wagener, and N.~Barker, ``Some different types of essential hypertension: their course and prognosis,'' \textit{Am. J. Med. Sci.}, vol. 197, no. 3, pp. 332--343, 1939.
\bibitem{wong2004} T.~Y. Wong and P.~Mitchell, ``Hypertensive retinopathy,'' \textit{New England J. Med.}, vol. 351, no. 22, pp. 2310--2317, 2004.
\bibitem{crawshaw2020} M.~Crawshaw, ``Multi-task learning with deep neural networks: a survey,'' \textit{arXiv:2009.09796}, 2020.
\bibitem{higgins2017} I.~Higgins et al., ``$\beta$-VAE: learning basic visual concepts with a constrained variational framework,'' in \textit{Proc. ICLR}, 2017.
\bibitem{hu2018} J.~Hu, L.~Shen, and G.~Sun, ``Squeeze-and-excitation networks,'' in \textit{Proc. CVPR}, 2018, pp. 7132--7141.
\bibitem{begoli2019} E.~Begoli, T.~Bhattacharya, and D.~Kusnezov, ``The need for uncertainty quantification in machine-assisted medical decision making,'' \textit{Nature Mach. Intell.}, vol. 1, pp. 20--23, 2019.
\bibitem{guo2017} C.~Guo et al., ``On calibration of modern neural networks,'' in \textit{Proc. ICML}, vol. 70, pp. 1321--1330, 2017.
\bibitem{leibig2017} C.~Leibig et al., ``Leveraging uncertainty information from deep neural networks for disease detection,'' \textit{Scientific Reports}, vol. 7, p. 17816, 2017.
\bibitem{selvaraju2017} R.~R. Selvaraju et al., ``Grad-CAM: visual explanations from deep networks via gradient-based localization,'' in \textit{Proc. ICCV}, 2017, pp. 618--626.
\bibitem{sayres2019} R.~Sayres et al., ``Using a deep learning algorithm and integrated gradients explanation to assist grading for diabetic retinopathy,'' \textit{Ophthalmology}, vol. 126, no. 4, pp. 552--564, 2019.
\bibitem{porwal2020} P.~Porwal et al., ``IDRiD: Diabetic retinopathy --- segmentation and grading challenge,'' \textit{Medical Image Analysis}, vol. 59, p. 101561, 2020.
\bibitem{staal2004} J.~J. Staal et al., ``Ridge-based vessel segmentation in color images of the retina,'' \textit{IEEE Trans. Med. Imag.}, vol. 23, no. 4, pp. 501--509, 2004.
\bibitem{loshchilov2019} I.~Loshchilov and F.~Hutter, ``Decoupled weight decay regularization,'' in \textit{Proc. ICLR}, 2019.
\end{thebibliography}

\end{document}
